\documentclass[11pt]{article}
\usepackage[a4paper,margin=22mm]{geometry}
\usepackage[no-math]{fontspec}
\setmainfont{Latin Modern Roman}
\usepackage{amsmath,amssymb,amsthm,xfrac,xspace,hyperref,graphicx,comment}
\excludecomment{DefTralics}
\providecommand{\C}{}
\providecommand{\bibliofont}{\small}
\newcommand{\ie}{i.e.,\xspace}
\newcommand{\eg}{e.g.,\xspace}
\newcommand{\etc}{etc.\xspace}
\newcommand{\cf}{cf.\xspace}
\newcommand{\resp}{resp.\xspace}
\newcommand{\smaller}{\small}
\newcommand{\Subsection}{\subsection}
\newcommand{\Subsubsection}{\subsubsection}
\newcommand{\skpt}{\smallskip}
\emergencystretch=3em
\input{author-preamble.tex}
\numberwithin{equation}{section}
\begin{document}
\begin{center}{\Large Stable item 1928}\end{center}
\paragraph{Source and attribution.}
Marc Briane, \emph{Homogenization of linear transport equations. A new approach}, Journal de l’École polytechnique — Mathématiques 7 (2020), publisher version, DOI 10.5802/jep.122. \href{https://jep.centre-mersenne.org/articles/10.5802/jep.122/}{Publisher article}. Authors' text: \href{https://creativecommons.org/licenses/by/4.0/}{CC BY 4.0}.
\paragraph{Editorial problem boundary.}
Prove only the first conclusion of Theorem 2.2 below using its uniform-divergence alternative (2.5). For the stated uniformly proper gradient rectification, positive bounded weights, directional strong convergence and uniformly bounded divergence, the selected weak weighted transport limit solves the printed effective equation with its stated initial data. The theorem\'s additional uniform-in-time strong \(L^2\) conclusion is not selected. Full notation, assumptions, weak-limit proof and cited prerequisite estimates are retained; surrounding article results are context only.
\paragraph{Difficulty (editorial): H2.}
Global rectifying diffeomorphisms and compensated compactness of gradient cofactors identify the effective transport flux. Composed test functions absorb the oscillatory velocities through strong convergence of the directional rectification coefficient, without periodicity or ergodicity.
\paragraph{Editorial transcription note.}
Original author language, mathematical content and ordering are unchanged. Publisher front matter is replaced by this independent layout wrapper. Original native statement, proof and macros remain editable; bibliography is recovered from the matching cached publisher HTML. No assistant-generated proof or mathematical repair is supplied.
\clearpage
\input{author-body.tex}
\begin{thebibliography}{99}
\bibitem{AHZ1} Y. Amirat, K. Hamdache \& A. Ziani - “Homogénéisation d’équations hyperboliques du premier ordre et application aux écoulements miscibles en milieu poreux” , Ann. Inst. H. Poincaré Anal. Non Linéaire 6 (1989) no. 5, p. 397-417
\bibitem{AHZ2} Y. Amirat, K. Hamdache \& A. Ziani - “Homogénéisation par décomposition en fréquences d’une équation de transport dans $\mathbf{R}^n$” , C. R. Acad. Sci. Paris Sér. I Math. 312 (1991) no. 1, p. 37-40
\bibitem{AHZ3} Y. Amirat, K. Hamdache \& A. Ziani - “Homogenisation of parametrised families of hyperbolic problems” , Proc. Roy. Soc. Edinburgh Sect. A 120 (1992) no. 3-4, p. 199-221
\bibitem{Bre} Y. Brenier - “Remarks on some linear hyperbolic equations with oscillatory coefficients” , in Third International Conference on Hyperbolic Problems, Vol. I, II (Uppsala, 1990) , Studentlitteratur, Lund, 1991, p. 119-130
\bibitem{Bri} M. Briane - “Isotropic realizability of fields and reconstruction of invariant measures under positivity properties. Asymptotics of the flow by a non-ergodic approach” , SIAM J. Appl. Dyn. Syst. 18 (2019) no. 4, p. 1846-1866
\bibitem{Cac} R. Caccioppoli - “Sugli elementi uniti delle trasformazioni funzionali: Un teorema di esistenza e di unicità ed alcune sue applicazioni” , Rend. Sem. Mat. Univ. Padova 3 (1932), p. 1-15
\bibitem{Dac} B. Dacorogna - Direct methods in the calculus of variations , Applied Math. Sci. , vol. 78 , Springer, New York, 2008
\bibitem{DiLi} R. J. DiPerna \& P.-L. Lions - “Ordinary differential equations, transport theory and Sobolev spaces” , Invent. math. 98 (1989) no. 3, p. 511-547
\bibitem{DEG} H. S. Dumas, J. A. Ellison \& F. Golse - “A mathematical theory of planar particle channeling in crystals” , Phys. D 146 (2000) no. 1-4, p. 341-366
\bibitem{DuGo} H. S. Dumas \& F. Golse - “The averaging method for perturbations of mixing flows” , Ergodic Theory Dynam. Systems 17 (1997) no. 6, p. 1339-1358
\bibitem{DGL} H. S. Dumas, F. Golse \& P. Lochak - “Multiphase averaging for generalized flows on manifolds” , Ergodic Theory Dynam. Systems 14 (1994) no. 1, p. 53-67
\bibitem{Gol1} F. Golse - “Moyennisation des champs de vecteurs et EDP” , in Journées “Équations aux Dérivées Partielles” (Saint Jean de Monts, 1990) , École Polytechnique, Palaiseau, 1990, Exp. no. XVI
\bibitem{Gol2} F. Golse - “Perturbations de systèmes dynamiques et moyennisation en vitesse des EDP” , C. R. Acad. Sci. Paris Sér. I Math. 314 (1992) no. 2, p. 115-120
\bibitem{HSD} M. W. Hirsch, S. Smale \& R. L. Devaney - Differential equations, dynamical systems, and an introduction to chaos , Elsevier/Academic Press, Amsterdam, 2013
\bibitem{HoXi} T. Y. Hou \& X. Xin - “Homogenization of linear transport equations with oscillatory vector fields” , SIAM J. Appl. Math. 52 (1992) no. 1, p. 34-45
\bibitem{Mul} S. Müller - “A surprising higher integrability property of mappings with positive determinant” , Bull. Amer. Math. Soc. (N.S.) 21 (1989) no. 2, p. 245-248
\bibitem{Mur} F. Murat - “Compacité par compensation” , Ann. Scuola Norm. Sup. Pisa Cl. Sci. (4) 5 (1978) no. 3, p. 489-507
\bibitem{Nis} Y. Nishimura - “Applications holomorphes injectives à jacobien constant de deux variables” , J. Math. Kyoto Univ. 26 (1986) no. 4, p. 697-709
\bibitem{ReSi} M. Reed \& B. Simon - Methods of modern mathematical physics, I. Functional analysis , Academic Press, Inc., New York, 1980
\bibitem{Sin} Y. G. Sinai - Introduction to ergodic theory , Princeton University Press, Princeton, NJ, 1976
\bibitem{Tar} L. Tartar - “Nonlocal effects induced by homogenization” , in Partial differential equations and the calculus of variations, Vol. II , Progr. Nonlinear Differential Equations Appl. , vol. 2 , Birkhäuser Boston, Boston, MA, 1989, p. 925-938
\bibitem{Tas} T. Tassa - “Homogenization of two-dimensional linear flows with integral invariance” , SIAM J. Appl. Math. 57 (1997) no. 5, p. 1390-1405
\end{thebibliography}
\end{document}
